\documentclass[a4paper,11pt]{article}
\usepackage[width=17cm,height=24cm]{geometry}
\usepackage[english]{babel}
\usepackage{amsmath,amssymb,mathrsfs,latexsym}

%%%%%%%%%% Start TeXmacs macros
\newcommand{\mathD}{\mathrm{D}}
\newcommand{\mathd}{\mathrm{d}}
\newcommand{\nobracket}{}
\newcommand{\tmaffiliation}[1]{\\ #1}
\newcommand{\tmem}[1]{{\em #1\/}}
\newcommand{\tmemail}[1]{\\ \textit{Email:} \texttt{#1}}
\newcommand{\tmop}[1]{\ensuremath{\operatorname{#1}}}
\newcommand{\tmscript}[1]{\text{\scriptsize{$#1$}}}
\newcommand{\tmtextit}[1]{{\itshape{#1}}}
\newenvironment{proof}{\noindent\textbf{Proof\ }}{\hspace*{\fill}$\Box$\medskip}
\newtheorem{definition}{Definition}
\newtheorem{lemma}{Lemma}
\newtheorem{theorem}{Theorem}
%%%%%%%%%% End TeXmacs macros

\newcommand{\cS}{\mathscr{S}}
\newcommand{\cN}{N}
\newcommand{\cA}{A}
\newcommand{\cQ}{Q}
\newcommand{\cC}{\mathscr{C}}
\newcommand{\cR}{R}
\newcommand{\cM}{M}
\newcommand{\FF}{\mathscr{F}}
\newcommand{\HH}{\mathscr{H}}
\newcommand{\cyl}{\mathcal{C} \tmop{yl}}
\newcommand{\JJ}{\mathcal{J}}
%

%

%


\begin{document}

\title{Energy solutions of KPZ are unique}

\author{
  M. Gubinelli
  \tmaffiliation{
  Hausdorff Center for Mathematics\\
   \& Institute for Applied Mathematics\\
   Universit{\"a}t Bonn}
%  \tmemail{}
  \and
  N. Perkowski
  \tmaffiliation{Humboldt-Universit{\"a}t zu Berlin\\
  Institut f{\"u}r Mathematik }
  \tmemail{perkowsk@math.hu-berlin.de}
}

\date{September 2015}

\maketitle

\begin{abstract}
Energy solutions is a weak notion of solution to the KPZ equation which was
introduced by Gon{\c c}alves and Jara [ARMA 212 (2013)] to describe the large
scale fluctuations of a wide class of weakly asymmetric particle systems. In
this paper we prove that energy solutions of the periodic, stationary,
stochastic Burgers equation are unique and in particular they coincide with
the Hopf--Cole solution. To obtain uniqueness we use the reformulation of the
energy solution concept given by Gubinelli and Jara [SPDE--AC 1 (2013)] and
an observation by Funaki and Quastel [ArXiv:1407.7310]. 
\end{abstract}

\section{Introduction}

Energy
solutions~{\cite{goncalves_nonlinear_2013,gubinelli_regularization_2013,gubinelli_lectures_2015}}
are a tool to give meaning to (stationary) solutions to the periodic
stochastic Burgers equation (SBE)
\begin{equation}
  \partial_t u = \Delta u + \partial_x u^2 + \partial_x \xi \label{eq:sbe}
\end{equation}
defined on the torus $\mathbb{T}=\mathbb{R}/ (2 \pi \mathbb{Z})$ with
initial condition $u_0 \sim \mu$ where $\mu$ is the white noise on $L^2_0
(\mathbb{T})$ (the space of square integrable functions on $\mathbb{T}$ with
zero mean), $\Delta$ is the periodic Laplacian and $\partial_t, \partial_x$
denote the time and space derivative respectively. The term $\xi$ is a
gaussian centered space--time white noise with covariance $\mathbb{E} [\xi
(t, x) \xi (s, y)] = \delta (t - s) \delta (x - y)$.

Energy solutions were introduced by Gon{\c c}alves and
Jara~{\cite{goncalves_nonlinear_2013}} in order to describe the \ fluctuations
of a large class of weakly asymmetric interacting particle system.
Subsequently Gubinelli and Jara~{\cite{gubinelli_regularization_2013}} used
the language of controlled paths to give a more abstract framework where
energy solutions are seen as some kind of ``weak solutions'' to the SBE. An
introductory exposition is present in the lecture
notes~{\cite{gubinelli_lectures_2015}} of Gubinelli and Perkowski.

The main difficulty in eq.~(\ref{eq:sbe}) is due to the interpretation of the
nonlinearity $\partial_x u^2$. Indeed any reasonable candidate solution
to~(\ref{eq:sbe}) features trajectories which, at fixed time and as functions
of the space variable, look like white noise on $\mathbb{T}$ so that the term
$\partial_x u^2$ does not have an obvious definition.

Bertini and Giacomin~{\cite{BertiniGiacomin1997}} determined the fluctuations
of the weakly asymmetric exclusion process (WASEP) by heavily exploiting the
fact that it is expected that any reasonable solution to~(\ref{eq:sbe}) should
be obtained as the Cole--Hopf transform $u = \partial_x \log \psi$ where
$\psi$ is the solution of the stochastic (linear, multiplicative) heat
equation (SHE):
\begin{equation}
  \mathd \psi = \Delta \psi + \psi \xi \label{eq:she}
\end{equation}
where $\xi$ is the space--time white noise and where the product $\psi \xi$
has to be understood as an Ito product between $\psi$ and a the
time--differential $\mathd W$ of a cylindrical Brownian motion on $L^2
(\mathbb{T})$~{\cite{prato_stochastic_2008}}. This last equation can be
solved via standard stochastic techniques giving a unique strong solution
$\psi$ which is space--time H{\"o}lder continuous and stricly positive for
positive times. In this way the relation $u = \partial \log \psi$ identifies a
unique generalized process which is called the Cole--Hopf solution of the
SBE~(\ref{eq:sbe}). Using a microscopic version of the Cole--Hopf transform
due to G{\"a}rtner, Bertini and Giacomin~{\cite{BertiniGiacomin1997}} were
able to prove that after rescaling and recentering the fluctuation field of
the WASEP converged to a process having the law of the Cole--Hopf solution.
This result relies on the precise structure of the WASEP and on the
availability of a Cole--Hopf correspondence at the microscopic level. However
it is expected (as {\tmem{the weak KPZ conjecture}} goes) that the fluctuation
field of a large class of microscopic non--equilibrium one--dimensional
particle system should likewise converge to the same object Hopf--Cole
solution. Gon{\c c}alves and Jara~{\cite{goncalves_nonlinear_2013}} were able
to prove the tightness of the fluctuation field of a large class of weakly
asymmetric interacting particle systems which do not enjoy a microscopic
Hopf--Cole transformation. The difficulty they encountered is that
caracterizing the limiting fluctuation field $u$ without the Hopf--Cole
correspondence hits the difficulty of describing the dynamics of $u$ in term
of~(\ref{eq:sbe}) which is not a well--posed equation as it stands, even
interpreted as a weak solution in the PDE sense. Their way out of this impasse
has been to give a new notion of solution which however, up to know, were not
know of enjoy uniqueness. In this way they were able to give a partial
caracterization of the accumulation points of the fluctuation field $u$
without being able to describe exactly the limit.

Note that in the meanwhile Hairer~{\cite{hairer_solving_2011}} has found a
notion of solution to the KPZ equation
\begin{equation}
  \partial_t h = \Delta h + (\partial_x h)^2 - \infty + \xi \label{eq:kpz}
\end{equation}
which enjoy a good existence and uniqueness theory and from which solution of
the SBE can be reobtained as $u = \partial_x h$, solving a long standing
problem of mathematical physics. The problem with Hairer's solution theory for
KPZ is that it is not simple to check that the limiting fluctuation fields of
the microscopic models satisfy the conditions of the theory. This difficulty
persists also in other recent (but, in this case, mostly equivalent) methods
of giving a meaning to the SBE: Hairer's regularity
structures~{\cite{Hairer2014Regularity}} or Gubinelli, Imkeller and
Perkowski's paracontrolled
distributions~{\cite{gubinelli_paracontrolled_2015,gubinelli_lectures_2015}}.

The aim of the present work is to show that stationary energy solutions
of~(\ref{eq:sbe}), as defined by Gubinelli and Jara
in~{\cite{gubinelli_regularization_2013}}, are unique and coincide with the
Hopf--Cole solution. Existence was already obtained by standard Galerkin
methods~{\cite{gubinelli_regularization_2013}}. In this way we fill a gap
which could lead to possibly estabilish the weak KPZ universality conjecture
in a quite large extent, at least in the stationary situation.

Our proof was motivated by the recent paper of Funaki and
Quastel~{\cite{funaki_kpz_2014}} in which they estabilish the invariance of
the white noise for the Hopf--Cole solution. In our framework invariance comes
for free from Galerkin approximations but their strategy can be adapted to
show that energy solution can be transformed via Hopf--Cole to give solutions
of the SHE. Uniqueness for the SHE then implies easily uniqueness for energy
solutions. In this way the full power of the weak formulation comes into play
since even if Hopf--Cole is not available at the microscopic level (or for
other kind of approximations, like Galerkin) the limiting energy solution does
allow the Hopf--Cole transformation. Note that our results does not seems to
apply to the original formulation of Gon{\c c}alves and Jara since it is not
clear how to perform in their setting the computations leading to the control
of the remainder in the Hopf--Cole transformation.

Extension of the uniqueness result in $\mathbb{R}$ is in progress and we do
not expect additional difficulties. Applications to the convergence toward KPZ
of the stationary version of the model considered by Hairer and
Quastel~{\cite{hairer-quastel}} are under investigation. \ \ \

\

\paragraph{Notation.}Let $\cS$ the space of $C^{\infty}$ functions on
$\mathbb{T}$ and by $\cS'$ the dual space of distributions. For $k \in
\mathbb{Z}$ we denote by $e_k (x) = \exp (- i k x) / (2 \pi)^{1 / 2}$ the
Fourier basis. For any $f \in \cS'$ let $\hat{f} = \FF f$ the Fourier
transform defined by $\FF f (k) = \langle e_k, f \rangle$. We have also $f (x)
= \sum_k \langle e_k, f \rangle e_k (x)$ and Plancherel formula $\langle f, g
\rangle_{L^2} = \sum_{k \in \mathbb{Z}} (\FF f (k))^{\ast} \FF g (k)$. Let
$\Pi_0$ the projection from $L^2 (\mathbb{T})$ to $L^2_0 (\mathbb{T})$ and
$\Pi^N_0$ the projection on the span of $(e_k)_{0 < | k | \leqslant N}$ and by
$\Pi^N$ the projection on the the span of $(e_k)_{0 \leqslant | k | \leqslant
N}$. Let $H^{\alpha}$ the standard Sobolev spaces on $\mathbb{T}$ of index
$\alpha \in \mathbb{R}$ and $H^{\alpha}_0 = \Pi_0 H^{\alpha}$ their
mean--zero counterparts.

The white noise $\mu$ can be realized on $H^{- \alpha}_0$ for some $\alpha > 1
/ 2$. If $u_0 \sim \mu$ then for all $k \in \mathbb{Z}\backslash \{ 0 \}$
$\FF u_0 (k)$ is a complex Gaussian variable such that $\FF u_0 (k)^{\ast} =
\FF u_0 (- k)$ and $\mathbb{E} \left| \FF u_0 (k) \right|^2 = 2$.

Denote $\cC_T V = C ([0, T], V)$ the space of continuous functions from $[0,
T]$ to the Banach space $V$ endowed with the supremum norm and with
$\cC^{\gamma}_T V = C^{\gamma} ([0, T], V)$ the subspace of
$\gamma$-H{\"o}lder continuous functions in $\cC_T V$ with the
$\gamma$-H{\"o}lder norm.

We introduce also a fixed family of smoothing operators indexed by $L
\geqslant 1$. Let $q : \mathbb{T} \rightarrow \mathbb{R}_+$ be a even smooth
function of compact support around $0$ and such that $\int_{\mathbb{T}} q (x)
\mathd x = 1$. Let $\delta^L (y) = Lq (L y)$, $\delta^L_x (y) = \delta^L (x -
y)$, $\rho^L_x = \delta^L_x - 1 / (2 \pi)$ and $\JJ^L_0 f = \rho^L_{} \ast f$.
We let also $\rho^{\infty}_x (y) = \delta_x (y) - 1 / (2 \pi)$ and then $\JJ_0
f = \rho^{\infty} \ast f = \Pi_0 f$.

\

\section{Controlled processes and energy solutions}

In this section we define stationary energy solutions of~(\ref{eq:sbe}),
following~{\cite{gubinelli_regularization_2013}}. We introduce first a class
of processes $u$ which at fixed time are distributed like the (zero--mean)
white noise on $\mathbb{T}$ but for which the nonlinear term $\partial_x u^2$
is defined as a space--time distribution. In this class of processes it makes
then sense to look for solutions of the SBE~(\ref{eq:sbe}).

\begin{definition}
  [Controlled process]\label{def:controlled}Denote with $\mathcal{Q}$ the
  space of pairs $(u, \mathcal{A})_{0 \leqslant t \leqslant T}$ of generalized
  stochastic processes with continuous paths in $\cS'$ such that
  \begin{itemize}
    \item[i)] the law of $u_t$ is the white noise $\mu$ for all $t \in [0,
    T]$;
    
    \item[ii)] For any test function $\varphi \in \cS$ the process $t \mapsto
    \mathcal{A}_t (\varphi)$ is a.s. of zero quadratic variation,
    $\mathcal{A}_0 (\varphi) = 0$ and the pair $(u (\varphi), \mathcal{A}
    (\varphi))_{0 \leqslant t \leqslant T} $ satisfies the equation
    \begin{equation}
      \label{eq:controlled-decomposition} u_t (\varphi) = u_0 (\varphi) +
      \int_0^t u_s (\Delta \varphi) \mathd s +\mathcal{A}_t (\varphi) +
      \partial_x M_t (\varphi)
    \end{equation}
    where $(M_t (\varphi))_{0 \leqslant t \leqslant T}$ is a martingale with
    respect to the filtration generated by $(u, \mathcal{A})_{0 \leqslant t
    \leqslant T}$ with quadratic variation $[M (\varphi)]_t = 2 t \| \varphi
    \|_{L^2 (\mathbb{T})}^2$;
    
    \item[iii)] the reversed processes $\hat{u}_t = u_{T - t}$,
    $\hat{\mathcal{A}}_t = - (\hat{\mathcal{A}}_T -\mathcal{A}_{T - t})$
    satisfies the same equation with respect to its own filtration.
  \end{itemize}
\end{definition}

When $\mathcal{A}= 0$ the process $(X, 0) \in \mathcal{Q}$ is the stationary
Ornstein--Uhlenbeck (OU) process with invariant measure $\mu$. It is the
unique solution of the SPDE
\[ \mathd X = \Delta X \mathd t + \partial_x \mathd M \]
with initial condition $u_0 \sim \mu$. Allowing $\mathcal{A} \neq 0$ has the
intuitive meaning of considering perturbation of the OU process with zero
quadratic variation antisymmetric drifts. In this sense we say that a couple
$(u, \mathcal{A}) \in \mathcal{Q}$ is a process {\tmem{controlled}} by the
Ornstein--Uhlenbeck process.

\

Controlled processes allow the definition of some interesting non--linear
functionals. We do not insists on generality here, all we need is that the
Burgers drift makes sense as a space--time distribution:

\begin{theorem}
  Assume that $(u, \mathcal{A}) \in \mathcal{Q}$ and \ for any $N \geqslant 1$
  and $0 \leqslant t \leqslant T$. Let
  \[ \mathcal{B}^N_t (\varphi) = \int_0^t \partial_x \left( \JJ^N_0 u_s
     \right)^2 (\varphi) \mathd s \]
  then $(\mathcal{B}^N_t)_{N \geqslant 1}$ converges in probability in $C^{1 /
  2 +} H^{- 1 -} \cap C^{0 +} H^{0 -}$, we denote the limit by
  \[ \int_0^t \partial_x u_s^2 (\varphi) \mathd s. \]
\end{theorem}

\begin{proof}
  The proof can be found
  in~{\cite{gubinelli_regularization_2013,gubinelli_lectures_2015}} for
  approximations given by $\tilde{\mathcal{B}}^N_t (\varphi) = \int_0^t
  \partial_x (\Pi^N_0 u_s)^2 (\varphi) \mathd s$. The estimates required to go
  from $\tilde{\mathcal{B}}^N_t$ to $\mathcal{B}^N_t$ simply uses the fact
  that $(s, x) \mapsto \JJ^N_0 u_s (x) - \Pi^N_0 u_s (x)$ is a smooth Gaussian
  random field which converges to zero as a continuous space-time function. 
\end{proof}

We are ready to give a rigorous meaning to the SBE~(\ref{eq:sbe}) in the class
of controlled processes:

\begin{definition}
  A controlled process $(u, \mathcal{A}) \in \mathcal{Q}$ is an energy
  solution of SBE iff
  \[ \mathcal{A}_t (\varphi) = \int_0^t \partial_x u_s^2 (\varphi) \mathd s.
  \]
  almost surely for all $t \in [0, T]$ and $\varphi \in \cS$.
\end{definition}

Existence of energy solutions of SBE can be proved via Galerkin approximations
or approximation of the non--linear term as done
in~{\cite{gubinelli_regularization_2013,gubinelli_lectures_2015}}. The aim of
this note is to show that energy solutions of SBE are in fact adapted to the
noise $M$ and pathwise unique. In order to state the result let $I u_0$ the
unique centered Gaussian random field such that $\partial_x I u_0 = u_0$ and
$\int_{\mathbb{T}} I u_0 (x) \mathd x = 0$.

\begin{theorem}
  \label{th:uniqueness}There exists only one energy solution $u$ of SBE and it
  is given by the Cole--Hopf transformation
  \[ u_t (x) = \partial_x \log \psi_t (x) \]
  where $\psi$ is the (It{\^o}) solution of SHE :
  \[ \mathd \psi_t = \Delta \psi_t + \psi_t \mathd M_t \]
  with the inital condition given by $\psi_0 (x) = \exp (I u_0 (x))$.
\end{theorem}

The rest of the paper is devoted to the proof of this theorem.

\section{Regularization by noise for controlled processes}

In this section $u$ will always denote a generic controlled process in
$\mathcal{Q}$, not necessarily an energy solution to SBE. Controlled processes
have regularization properties coming from the fast decorrelation in time of
the OU process associated to the Laplacian $\Delta$. Using an It{\^o} formula
and following the classic approach of Kipnis--Varadhan, we can replace
time--averages of some functional with a sum of forward and backward
martingales whose quadratic variation is controlled in terms of a more regular
(or smaller) functional.

For any test function $\varphi \in \cS$ the processes $(u_t (\varphi))_{t \in
[0, T]}$ and $(\hat{u}_t (\varphi))_{t \in [0, T]}$ are Dirichlet processes:
the sum of a martingale and a zero quadratic variation process. This is
compatible with the regularity of our solutions and there is no clue that
solutions of SBE are distributional semimartingales. Dirichlet processes
(wheter Markov or not) can be handled via the stochastic calculus via
regularization developed by Russo and Vallois~{\cite{russo_elements_2007}}. In
this approach the It{\^o} formula holds also for Dirichlet processes and if $X
= (X^i)_{i = 1, \ldots, k}$ is an $\mathbb{R}^k$ valued Dirichlet process and
$g$ is a $C^2 (\mathbb{R}^k ; \mathbb{R})$ function then we have
\begin{equation}
  g (X_t) = g (X_0) + \sum_{i = 1}^k \int_0^t \partial_i g (X_s) \mathd^-
  X^i_s + \frac{1}{2}  \sum_{i, j = 1}^k \int_0^t \partial^2_{i, j} g (X_s)
  \mathd^- [X^i, X^j]_s \label{eq:ito-fv}
\end{equation}
where $\mathd^-$ denotes the forward integral and $[X, X]$ the quadratic
covariation of the vector process $X$. Decomposing $X = M + N$ as the sum of a
martingale $M$ and a zero quadratic variation process $N$ we have $[X, X] =
[M, M]$ and
\[ g (X_t) = g (X_0) + \sum_{i = 1}^k \int_0^t \partial_i g (X_s) \mathd^-
   M^i_s + \sum_{i = 1}^k \int_0^t \partial_i g (X_s) \mathd^- N^i_s \]
\[ + \sum_{i, j = 1}^k \frac{1}{2}  \int_0^t \partial^2_{i, j} g (X_s)
   \mathd^- [M^i, M^j]_s \]
where now $\mathd^- M$ coincide with the usual It{\^o} integral and $[M, M]$
is the usual quadratic variation of the martingale $M$. The integral $\int_0^t
\partial_i g (X_s) \mathd^- N^i_s$ is well-defined due to the fact that all
the other terms in this formula are well defined. The case the function $g$
depends explicitly on time can be handled by the above formula by considering
time as an additional ($0$-th) component of the process $X$ and using the fact
that $[X^i, X^0] = 0$ for all $i = 1, .., k$. In the computations which
follows we will only need to apply the It{\^o} formula to smooth functions.

\

On $H = L^2_0 (\mathbb{T})$ we consider the space of smooth cylinder
functions $\cyl$ which depends only on finitely many coordinates on the basis
$\{e_k \}_{k \in \mathbb{Z}_0}$ and we consider the Malliavin derivative $D :
L^2 (\mu) \to L^2 (\mu) \otimes H$ defined as $D \varphi (x) = \sum_{k \in
\mathbb{Z}_0} \partial_{x_k} \varphi (x) e_k$ for $\varphi \in \cyl$ where
$x_k = \langle e_k, x \rangle$ are the coordinates of $x$. We let $D_k \varphi
= \langle e_k, D \varphi \rangle$.

\

With this notations we can write the generator $L_0$ of the
Ornstein-Uhlenbeck process associated to the operator $\Delta$ as
\begin{equation}
  \label{eq:generator-ou} L_0 \varphi (x) = \sum_{| k | > 0} |k|^2  \left( -
  x_k D_k \varphi (x) + \tfrac{1}{2} D_{- k} D_k \varphi (x) \right) .
\end{equation}


For any smooth cylinder function $h : [0, T] \times \JJ_N H \to \mathbb{R}$
the It{\^o} formula for the finite quadratic variation process $(u^N_t =
\Pi^N_0 u_t)_t$ gives
\[ h (t, u^N_t) = h (0, u^N_0) + \int_0^t (\partial_s + L^N_0) h (s, u^N_s)
   \mathd s + \int_0^t Dh (s, u^N_s) \mathd \JJ^N \mathcal{A}_s + M^+_t \]
where
\[ L^N_0 h (s, x) = \sum_{0 < |k| \le N} |k|^2  (- x_k D_k h (s, x) +
   \frac{1}{2} D_k D_{- k} h (s, x)) \]
is the restriction of the operator $L_0$ to $\Pi^N H$ and where the martingale
part denoted $M^+$ has quadratic variation given by $[M^+]_t = \int_0^t
\mathcal{E}^N (h (s, \cdot)) (u^N_s) \mathd s$, where
\[ \mathcal{E}^N (\varphi) (x) = \frac{1}{2}  \sum_{0 < |k| \le N} |k|^2  |D_k
   \varphi (x) |^2 = \left\| \JJ_0^N D \varphi \right\|_{H^1_0}^2, \]
is the density of the Dirichlet form associated to the generator $L_0^N$. The
It{\^o} formula on the backward process reads
\[ h (T - t, u^N_{T - t}) = h (T, u^N_T) + \int_0^t (- \partial_s + L^N_0) h
   (T - s, u^N_{T - s}) \mathd s \]
\[ - \int_0^t Dh (T - s, u^N_{T - s}) \mathd \JJ^N \mathcal{A}_{T - s} + M^-_t
\]
with $[M^-]_t = \int_0^t \mathcal{E}^N (h (T - s, \cdot)) (u^N_{T - s}) \mathd
s$. Adding the two It{\^o} formulas we have the key equality
\begin{equation}
  \label{eq:key-representation} \int_0^t 2 L_0^N h (s, u^N_s) \mathd s = -
  M^+_t + M^-_{T - t} - M^-_T .
\end{equation}
which allows us to represent the time integral of $2 L_0^N h$ as a sum of
martingales which allows better control since their quadratic variation
depends only on $\mathcal{E}^N (h)$. From this we can prove easily
(see~{\cite{gubinelli_regularization_2013}}) the following lemma.

\begin{lemma}
  [It{\^o} trick]\label{lemma:ito-trick}Let $h : [0, T] \times \Pi_N H \to
  \mathbb{R}$ be a cylinder function. Then for any $p \geq 1$,
  \begin{equation}
    \label{eq:ito-trick} \left\| \sup_{t \in [0, T]} \left| \int_0^t L_0 h (s,
    u_s) \mathd s \right| \right\|_{L^p (\mathbb{P}_{\mu})} \lesssim_p T^{1 /
    2} \sup_{s \in [0, T]} \| \mathcal{E}(h (s, \cdot)) \nobracket \|^{1 /
    2}_{L^{p / 2} (\mu)}
  \end{equation}
  where $\mathcal{E} (\varphi) (x) = \frac{1}{2}  \sum_{| k | > 0} |k|^2  |D_k
  \varphi (x) |^2$. In the particular case $h (s, x) = e^{a (T - s)} \tilde{h}
  (x)$ for some $a \in \mathbb{R}$ we have the improved estimate
  \begin{equation}
    \left\| \int_0^T e^{a (T - s)} L_0  \tilde{h} (u_s) \mathd s \right\|_{L^p
    (\mathbb{P}_{\mu})} \lesssim_p \left( \frac{1 - e^{2 aT}}{2 a} \right)^{1
    / 2} \| \mathcal{E}(\tilde{h}) \|^{1 / 2}_{L^{p / 2} (\mu)} .
    \label{eq:ito-trick-conv}
  \end{equation}
\end{lemma}

If we take $p = 2$ in Lemma~\ref{lemma:ito-trick} we have
\[ \mathbb{E} \left[ \sup_{t \in [0, T]} \left| \int_0^t L_0 h (u_s) \mathd s
   \right|^2_{} \right] \lesssim T \| h \|_1^2 \]
where
\[ \| h \|_1^2 =\mathbb{E}\mathcal{E} (h (u_0)) =\mathbb{E} \| \mathD h
   (u_0) \|_{H^1}^2 \]
is the Dirichlet form of the OU process.

\

We denote by $\HH^1$ the completion of the space of smooth cylinder functions
with respect to the norm $\| \ast \|_1$ and denote with $\HH^{- 1}$ the dual
of $\HH^1$ with norm
\[ \| f \|_{- 1} = \sup_{g \in L^2 (\mu)} (2 \langle f, g \rangle - \| g
   \|^2_1) \]
(see~{\cite{fluctuations_2012}}, Chap.~2 for details).

\

In order to exploit the It{\^o} trick of Lemma~\ref{lemma:ito-trick} we need
to be able to solve the Poisson equation $L_0 h_f = f$ for given function $f$.
Sometimes we can do it explicitly and estimate directly the Dirichlet energy
of the solution. In other situations however it is preferable to have a bound
involving directly $f$ without need to find $h_f$ explicitly. (In order not to
bother us with domain considerations the reader can think of cylindrical and
smooth $f, h$ since this will be enough for our purposes.) Fix an arbitrary
$h$ and consider the decomposition
\[ f = L_0 h + L_0 h - f \]
then
\[ \left\| \sup_{t \in [0, T]} \left| \int_0^t f (u_s) \mathd s \right|
   \right\|_{L^2 (\mathbb{P})} \leqslant \left\| \sup_{t \in [0, T]} \left|
   \int_0^t L_0 h (u_s) \mathd s \right| \right\|_{L^2 (\mathbb{P})} +
   \left\| \sup_{t \in [0, T]} \left| \int_0^t (L_0 h - f) (u_s) \mathd s
   \right| \right\|_{L^2 (\mathbb{P})} \text{} \]
\[ \leqslant \left\| \sup_{t \in [0, T]} \left| \int_0^t L_0 h (u_s) \mathd s
   \right| \right\|_{L^2 (\mathbb{P})} + \int_0^T \| (L_0 h - f) (u_s)
   \|_{L^2 (\mathbb{P})} \mathd s \]
\[ \leqslant C T^{1 / 2}  \| h \|_1 + T \| L_0 h - f \|_{L^2 (\mu)} \]
Now note that, for fixed $\lambda > 0$ a solution $h_{\lambda}$ to $(\lambda -
L_0) h_{\lambda} = - f$ is given explicitly by
\[ h_{\lambda} (x) = - \int_0^{\infty} e^{- \lambda s} (P_s^{\tmop{OU}} f) (x)
   \mathd s \]
where $P^{\tmop{OU}}$ is the contraction semigroup generated by $L_0$ on $L^2
(\mu)$. Then
\[ \langle h_{\lambda}, - f \rangle = \langle h_{\lambda}, (\lambda - L_0)
   h_{\lambda}, \rangle = \lambda \| h_{\lambda} \|^2_{L^2} + \| h_{\lambda}
   \|_1^2 \]
and if $f \in \HH^{- 1}$ we have
\[ \lambda \| h_{\lambda} \|^2_{L^2} + \| h_{\lambda} \|_1^2 = | \langle
   h_{\lambda}, f \rangle | \leqslant \| f \|_{- 1} \| h_{\lambda} \|_1 \]
which implies $\| h_{\lambda} \|_1 \leqslant \| f \|_{- 1}$ and $\|
h_{\lambda} \|_{L^2} \leqslant \lambda^{- 1 / 2} \| f \|_{- 1}$. So noting
that $L_0 h_{\lambda} - f = \lambda h_{\lambda}$ we have


\[ C T^{1 / 2}  \| h_{\lambda} \|_1 + T \| L_0 h_{\lambda} - f \|_{L^2 (\mu)}
   = C T^{1 / 2}  \| h_{\lambda} \|_1 + T \| \lambda h_{\lambda} \|_{L^2
   (\mu)} \leqslant C T^{1 / 2} \| f \|_{- 1} + T \lambda^{1 / 2} \| f \|_{-
   1} . \]
Taking $\lambda \rightarrow 0$ we end up with the following lemma which
extends the Kipnis--Varadhan lemma (cfr. Lemma 2.4
in~{\cite{fluctuations_2012}}) \ to controlled processes :

\begin{lemma}
  \label{lemma:kv}Assume that $f \in \HH^{- 1}$ then for every controlled
  process $u$ we have
  \[ \left\| \sup_{t \in [0, T]} \left| \int_0^t f (u_s) \mathd s \right|
     \right\|_{L^2 (\mathbb{P})} \leqslant C T^{1 / 2} \| f \|_{- 1} . \]
  with some constant $C$ which does not depends on $u$.
\end{lemma}

\section{Mapping to the SHE}

The strategy to prove Theorem~\ref{th:uniqueness} is the following. Fix $L
\geqslant 1$ and consider the process $\phi_t^L (x) = \exp (I u_t^L (x))$
where $u_t^L = \JJ^L_0 u_t$ and $I = \partial_x^{- 1}$ is the antiderivative
defined on $L^2_0 (\mathbb{T})$ by $\partial_x I \varphi = \varphi$ for all
$\varphi \in L^2_0 (\mathbb{T})$. Let $I^{\ast} = - I$ the adjoint of $I$ and
$\rho_x^L = \JJ^L_0 \delta_x$ is a smooth test function of zero mean. Now $I
\rho^L_x (y) = (\delta^L \ast \Theta) (y - x)$ where $\Theta : \mathbb{T}
\rightarrow R$ is the odd periodic function given by $\Theta (y) = \frac{1}{2}
\tmop{sgn} (y) - y / (2 \pi)$ for $y \in (- \pi, \pi]$ with $\Theta (0) = 0$.

The process $(\phi_t^L (x))_{t, x}$ is smooth in $x$ and of finite quadratic
variation in $t$. Note that $\phi_t^L (x) = \exp (I u_t^L (x)) = \exp (u_t
(I^{\ast} \rho^L_x))$ is not a cylindrical function but by a direct
approximation procedure we can pretend it is and since $(u_t, \mathcal{A})$ is
controlled we have
\[ \mathd (I u_t^L) (x) = \mathd u_t (I^{\ast} \rho^L_x) = u_t (\Delta
   I^{\ast} \rho^L_x) \mathd t + \mathd \mathcal{A}_t (I^{\ast} \rho^L_x) -
   \mathd M_t (\partial_x I^{\ast} \rho^L_x) . \]
The It{\^o} formula for Dirichlet processes~(\ref{eq:ito-fv}) gives
\[ \mathd \phi_t^L (x) = \phi_t^L (x) (u_t (\Delta I^{\ast} \rho^L_x) \mathd t
   + \mathd \mathcal{A}_t (I^{\ast} \rho^L_x) - \mathd M_t (\partial_x
   I^{\ast} \rho^L_x)) + \phi_t (x) \langle \partial_x I^{\ast} \rho^L_x,
   \partial_x I^{\ast} \rho^L_x \rangle \mathd t \]
\[ = \phi_t^L (x) (u_t (\Delta I^{\ast} \rho^L_x) \mathd t + \mathd
   \mathcal{A}_t (I^{\ast} \rho^L_x) + \mathd M_t (\rho^L_x)) + \phi_t (x)
   \langle \rho^L_x, \rho^L_x \rangle \mathd t \]
Moreover since $\rho^L_x (y) = \rho^L_0 (y - x)$ we have $\Delta (I^{\ast}
\rho^L_x) = \Delta_x (I^{\ast} \rho^L_x)$ and
\[ \phi_t^L (x) u_t (\Delta I^{\ast} \rho^L_x) = \phi_t^L (x) \Delta_x u_t
   (I^{\ast} \rho^L_x) = \Delta_x \phi_t^L (x) - \phi_t^L (x) (\partial_x u_t
   (I^{\ast} \rho^L_x))^2 \]
\[ = \Delta_x \phi_t^L (x) - \phi_t^L (x) (u_t (\partial_x I^{\ast}
   \rho^L_x))^2 = \Delta_x \phi_t^L (x) - \phi_t^L (x) (u_t (\rho^L_x))^2 . \]
Being $\phi_t^L$ is smooth in space we can rewrite this as an approximate
stochastic heat equation (SHE)
\[ \mathd \phi_t^L (x) = \Delta_x \phi_t^L (x) \mathd t + \phi_t^L (x) (K^L -
   Q_t^L) \mathd t + \phi_t^L (x) \mathd M_t (\rho_x^L) + \mathd R_t^L (x) \]
where the remainder $R^L$ is given by
\[ R_t^L (x) = \int_0^t \phi_s^L (x) (\mathd \mathcal{A}_s (I^{\ast} \rho^L_x)
   - \JJ_0 (u_s (\rho^L_{\ast}))^2 (x) \mathd s - K^L \mathd s) \]
and where we introduced a constant $K^L$ and the one--dimensional process
\[ Q_t^L = \frac{1}{2 \pi} \int_{\mathbb{T}} [(u_t (\rho^L_x))^2 - \langle
   \rho^L_x, \rho^L_x \rangle] \mathd x \]
By definition of energy solution of SBE we have also
\begin{equation}
  \mathcal{A}_t (I^{\ast} \rho_x^L) = - \lim_{N \rightarrow 0} \int_0^t \left(
  \JJ^N_0 u_s \right)^2 (\partial_x I^{\ast} \rho^L_x) \mathd s = \lim_{N
  \rightarrow 0} \int_0^t \left( \JJ^N_0 u_s \right)^2 (\rho^L_x) \mathd s
  \label{eq:A-limit}
\end{equation}
Thanks to eq.~(\ref{eq:A-limit}) we can define for any $N \geqslant 1$ the
process
\[ R_t^{L, N} (x) = \int_0^t \phi_s^L (x) (\left( \JJ^N_0 u_s \right)^2
   (\rho^L_x) - \JJ_0 \left( (\JJ^L_0 u_s)^2 \right) (x) - K^{L, N}) \mathd s
\]
\[ = \int_0^t \phi_s^L (x) (\JJ^L_0 \left( \JJ^N_0 u_s \right)^2 (x) - \JJ_0
   \left( (\JJ^L_0 u_s)^2 \right) (x) - K^{L, N}) \mathd s \]
and observe that $\lim_{N \rightarrow \infty} R^{L, N} = R^L$ if we choose the
sequence of constants $K^{L, N}$ so that $K^{L, N} \rightarrow K^L$. In
Sec.~\ref{sec:remainder} below we prove the following lemma which is the key
for the convergence

\begin{lemma}
  \label{lemma:key}There exists a choice of $K^{L, N}$ such that for any
  $\varphi \in \cS$ we have $\mathbb{E} [R^{L, N}_t (\varphi)] = 0$ and
  \[ \lim_{L \rightarrow \infty} \lim_{N \rightarrow \infty} \mathbb{E}
     [\sup_{0 \leqslant t \leqslant T} (R^{L, N}_t (\varphi))^2] = 0 \]
  where $R^{L, N}_t (\varphi) = \langle R^{L, N}_t, \varphi \rangle_{L^2
  (\mathbb{T})}$. Moreover $K^{L, N} \rightarrow K^L$ as $N \rightarrow +
  \infty$ and $K^L \rightarrow K = - 1 / 12$ as $L \rightarrow + \infty$.
\end{lemma}

We are now ready to prove our main result.

\begin{proof}
  (of Theorem~\ref{th:uniqueness}) We have
  \[ R^L_t (\varphi) = \phi_t^L (\varphi) - \phi_0^L (\varphi) - \int_0^t
     \phi_s^L (\Delta \varphi) \mathd s - \int_0^t \phi_s^L (\varphi) (K^L -
     Q_s^L) \mathd s - \int_0^t \int_{\mathbb{T}} \phi_s^L (x) \varphi (x)
     \mathd M_s (\rho_x^L) \mathd x \]
  We already know that the family $(\phi_t^L (x))_{t, x}$ convergences in the
  space of continuous paths on $[0, T] \times \mathbb{T}$ to $\phi_t (x) =
  \exp ((I u_t) (x))$. It is not difficult also to show that the process
  $(Q_s^L)_s$ converges in $L^2 ([0, T])$ to a limit which we call $Q$. The
  quadratic variation of the martingale part converges also, so we have that,
  for fixed $\varphi \in \cS$
  \[ R_t (\varphi) = \lim_{L \rightarrow \infty} R^L_t (\varphi) \]
  \[ = \phi_t (\varphi) - \phi_0 (\varphi) - \int_0^t \phi_s (\Delta \varphi)
     \mathd s - \int_0^t \phi_s (\varphi) (K - Q_s) \mathd s - \int_0^t
     \int_{\mathbb{T}} \phi_s (x) \varphi (x) \mathd \left( \JJ_0 M_s \right)
     (x) \mathd x \]
  as a continuous process in $t \in [0, T]$. But from Lemma~\ref{lemma:key} we
  have also by Fatou
  \[ \mathbb{E} [\sup_{0 \leqslant t \leqslant T} (R^{}_t (\varphi))^2]
     \leqslant \liminf_{L \rightarrow \infty} \mathbb{E} [\sup_{0 \leqslant t
     \leqslant T} (R^L_t (\varphi))^2] \leqslant \liminf_{L \rightarrow
     \infty} \liminf_{N \rightarrow \infty} \mathbb{E} [\sup_{0 \leqslant t
     \leqslant T} (R^{L, N}_t (\varphi))^2] = 0 \]
  so we conclude that $\phi$ satisfy
  \[ \phi_t (\varphi) - \phi_0 (\varphi) - \int_0^t \phi_s (\Delta \varphi)
     \mathd s - \int_0^t \phi_s (\varphi) (K - Q_s) \mathd s - \int_0^t
     \int_{\mathbb{T}} \phi_s (x) \varphi (x) \mathd \left( \JJ_0 M_s \right)
     (x) \mathd x = 0 \]
  almost surely for all $\varphi \in \cS$ and $t \in [0, T]$. Note also that
  the process
  \[ \tilde{\phi}_t (x) = \phi_t (x) \exp \left( M_t (1) - \frac{t}{2} + Kt -
     \int_0^t Q_s \mathd s \right) = \exp \left( (I u_t) (x) + M_t (1) -
     \frac{t}{2} + Kt - \int_0^t Q_s \mathd s \right) \]
  satisfies
  \[ \tilde{\phi}_t (\varphi) - \tilde{\phi}_0 (\varphi) - \int_0^t
     \tilde{\phi}_s (\Delta \varphi) \mathd s - \int_0^t \int_{\mathbb{T}}
     \tilde{\phi}_s (x) \varphi (x) \mathd M_s (x) \mathd x = 0. \]
  However we also know that there exists a unique strictly positive solution
  $\psi$ for this equation which is adapted to the filtration generated by
  $M$, so we must have $\tilde{\phi}_t (x) = \psi_t (x)$ and in particular
  \[ (I u_t) (x) + M_t (1) - \frac{t}{2} + Kt - \int_0^t Q_s \mathd s = \log
     \tilde{\phi}_t (x) = \log \psi_t (x) \]
  By differentiating this equality (in the sense of distributions) and
  projecting away the constants we get
  \[ u_t (x) = \partial_x \log \psi_t (x) \]
  which shows that the only energy solution of the SBE is the Cole--Hopf
  solution obtained from solving the SHE associated to the noise $M$.
\end{proof}

\section{Convergence of the remainder}\label{sec:remainder}

Let
\[ r^{L, N} (u_s, x) = \phi_s^L (x) (\JJ^L_0 \left( \JJ^N_0 u_s \right)^2 (x)
   - \JJ_0 (\JJ^L_0 u_s)^2 (x) - K^{L, N}) \]
so that $R_t^{L, N} (x) = \int_0^t r^{L, N} (u_s, x) \mathd s$. Using
Lemma~\ref{lemma:kv} we can estimate
\[ \mathbb{E} [\sup_{0 \leqslant t \leqslant T} (R^{L, N}_t (\varphi))^2]
   \lesssim T \| r^{L, N} (\cdot, \varphi) \|_{- 1}^2 \]
where
\[ \| r^{L, N} (\cdot, \varphi) \|_{- 1}^2 = \sup_{\Phi \in L^2 (\mu)}
   2\mathbb{E} [r^{L, N} (u_0, \varphi) \Phi (u_0)] - \| \Phi \|_1^2 \]
with $\| \Phi \|_1^2 =\mathbb{E} \| \mathD \Phi \|_{H^1}^2$ in terms of the
Malliavin derivative $D$ associated to the measure $\mu$. We prove below that
we can choose $K^{L, N}$ so that $\mathbb{E} [r^{L, N} (u_t, x)] =\mathbb{E}
[r^{L, N} (u_0, x)] = 0$ for all $x \in \mathbb{T}$. This is necessary in
order for $\| r^{L, N} (\cdot, \varphi) \|_{- 1}$ to be finite for all
$\varphi$. At this point everything boils down to control
\[ \mathbb{E} [r^{L, N} (u_0, \varphi) \Phi (u_0)] \]
and show that it goes to zero as $N \rightarrow + \infty$ and $L \rightarrow +
\infty$ where $\varphi$ is a smooth test function and $\Phi \in L^2 (\mu)$ is
such that $\| \Phi \|_1 < + \infty$.

\

For any $n \geqslant 1$ we denote by
\[ W_n (h) = \int_{\mathbb{T}^n} h (z_1, \ldots, z_n) W (\mathd z_1 \cdots
   \mathd z_n) \]
the $n$-th chaos of the white noise $u_0$. \ It is also possible to check that
when $\Psi$ is in the domain of the Malliavin derivative $D$ we can partially
integrate by parts in the expression $\mathbb{E} [W_n (h) \Psi]$ to obtain
the equality
\[ \mathbb{E} [W_n (h) \Psi] = \int_{\mathbb{T}} \mathbb{E} [W_{n - 1} (h
   (z_1, \cdot, \ldots, \cdot)) D_{z_1} \Psi] \mathd z_1 \]
this will be used repeatedly below.

\

Observe that the random variable $(\JJ^L_0 \left( \JJ^N_0 u_0 \right)^2 (x) -
\JJ_0 \left( (\JJ^L_0 u_0)^2 \right) (x))$ is an element of the second chaos
of $u_0$. Let us compute its kernel:
\[ \JJ^L_0 \left( \JJ^N_0 u_0 \right)^2 (x) = \int_{\mathbb{T}^2} \left[
   \int_{\mathbb{T}} \mathd y \rho_x^L (y) \rho^N_y (z_1) \rho^N_y (z_2)
   \right] W (\mathd z_1 \mathd z_2) \]
and
\[ \JJ_0 \left( (\JJ^L_0 u_0)^2 \right) (x) = \int_{\mathbb{T}^2} \left[
   \int_{\mathbb{T}} \mathd y \rho_x^{\infty} (y) \rho^L_y (z_1) \rho^L_y
   (z_2) \right] W (\mathd z_1 \mathd z_2) \]
where
\[ \rho_x^{\infty} (y) = \lim_{K \rightarrow \infty} \rho^K_y (y) = \delta_x
   (y) - \frac{1}{2 \pi} \]
is the distributional kernel of $\JJ_0$. Let
\[ g_x^{L, N} (z_1, z_2) = \int_{\mathbb{T}} \mathd y (\rho_x^L (y) \rho^N_y
   (z_1) \rho^N_y (z_2) - \rho_x^{\infty} (y) \rho^L_y (z_1) \rho^L_y (z_2))
\]
so that
\[ (\JJ^L_0 \left( \JJ^N_0 u_0 \right)^2 (x) - \JJ_0 \left( (\JJ^L_0 u_0)^2
   \right) (x)) = W_2 (g_x^{L, N}) = \int_{\mathbb{T}^2} g_x^{L, N} (z_1,
   z_2) W (\mathd z_1 \mathd z_2) . \]
We let also
\[ W_1 (g_x^{L, N} (z_1, \cdot)) = \int_{\mathbb{T}} g_x^{L, N} (z_1, z_2) W
   (\mathd z_2) . \]
Consider now
\[ \mathbb{E} [r^{\varepsilon} (u_0, \varphi) \Phi (u_0)] =
   \int_{\mathbb{T}} \varphi (x) \mathbb{E} [(W_2 (g_x^{L, N}) - K^{L, N})
   \phi_0^L (x) \Phi (u_0)] \mathd x \]
Partially integrating by parts $W_2$ we have
\[ \mathbb{E} [W_2 (g_x^{L, N}) \phi_0^L (x) \Phi (u_0)] = \int_{\mathbb{T}}
   \mathbb{E} [W_1 (g_x^{L, N} (z_1, \cdot)) \mathD_{z_1} [\phi_0^L (x) \Phi
   (u_0)]] \mathd z_1 \]
\[ = \int_{\mathbb{T}} \mathbb{E} [W_1 (g_x^{L, N} (z_1, \cdot)) \phi_0^L
   (x) \mathD_{z_1} \Phi (u_0)] \mathd z_1 + \int_{\mathbb{T}} \mathbb{E}
   [W_1 (g_x^{L, N} (z_1, \cdot)) \mathD_{z_1} \phi_0^L (x) \Phi (u_0)] \mathd
   z_1 \]
The second term can be again integrated by parts to obtain
\[ \int_{\mathbb{T}} \mathbb{E} [W_1 (g_x^{L, N} (z_1, \cdot)) \mathD_{z_1}
   \phi_0^L (x) \Phi (u_0)] \mathd z_1 = \int_{\mathbb{T}^2} g_x^{L, N} (z_1,
   z_2) \mathbb{E} [\mathD_{z_1, z_2}^2 \phi_0^L (x) \Phi (u_0)] \mathd z_1
   \mathd z_2 \]
\[ + \int_{\mathbb{T}^2} g_x^{L, N} (z_1, z_2) \mathbb{E} [\mathD_{z_1}
   \phi_0^L (x) \mathD_{z_2} \Phi (u_0)] \mathd z_1 \mathd z_2 . \]
while the first can be written as
\[ \int_{\mathbb{T}} \mathbb{E} [W_1 (g_x^{L, N} (z_1, \cdot)) \phi_0^L (x)
   \mathD_{z_1} \Phi (u_0)] \mathd z_1 = \int_{\mathbb{T}} \mathbb{E} [(W_1
   (g_x^{L, N} (z_1, \cdot)) \diamond \phi_0^L (x)) \mathD_{z_1} \Phi (u_0)]
   \mathd z_1 \]
\[ + \int_{\mathbb{T}^2} g_x^{L, N} (z_1, z_2) \mathbb{E} [\mathD_{z_2}
   \phi_0^L (x) \mathD_{z_1} \Phi (u_0)] \mathd z_1 \mathd z_2 . \]
where
\[ (W_1 (g_x^{L, N} (z_1, \cdot)) \diamond \phi_0^L (x)) = W_1 (g_x^{L, N}
   (z_1, \cdot)) \phi_0^L (x) - \int_{\mathbb{T}} g_x^{L, N} (z_1, z_2)
   \mathD_{z_2} \phi_0^L (x) \mathd z_2 \]
is a (partial) Wick contraction of the two terms. In the end we have the
decomposition
\[ \mathbb{E} [r^{L, N} (u_0, \varphi) \Phi (u_0)] = A^{L, N} + B^{L, N} +
   C^{L, N} \]
where
\[ A^{L, N} = \int_{\mathbb{T}} \varphi (x) \int_{\mathbb{T}} \mathbb{E}
   [(W_1 (g_x^{L, N} (z_1, \cdot)) \diamond \phi_0^L (x)) \mathD_{z_1} \Phi
   (u_0)] \mathd z_1 \mathd x \]
\[ B^{L, N} = \int_{\mathbb{T}} \varphi (x) \left[ \int_{\mathbb{T}^2}
   g_x^{L, N} (z_1, z_2) \mathbb{E} \left[ (\mathD_{z_1, z_2}^2 \phi_0^L (x))
   \Phi \left( u_0 \right) \right] \mathd z_1 \mathd z_2 - K^{L, N}
   \mathbb{E} [^{} \phi_0^L (x) \Phi (u_0)] \right] \mathd x \]
and
\[ C^{L, N} = 2 \int_{\mathbb{T}} \varphi (x) \int_{\mathbb{T}^2} g_x^{L, N}
   (z_1, z_2) \mathbb{E} [(\mathD_{z_1} \phi_0^L (x)) \mathD_{z_2} \Phi
   (u_0)] \mathd z_1 \mathd z_2 \mathd x. \]
\[ \  \]
So we can bound these three terms independently. In order to proceed, observe
that
\[ \mathD_{z_1} \phi_0^L (x) = \phi_0^L (x) (I^{\ast} \rho_x^L) (z_1) \]
and
\[ \mathD_{z_1, z_2}^2 \phi_0^L (x) = \phi_0^L (x) (I^{\ast} \rho_x^L) (z_1)
   (I^{\ast} \rho_x^L) (z_2) \]
so
\[ B^{L, N} = \int_{\mathbb{T}} \varphi (x) \left[ \int_{\mathbb{T}^2}
   g_x^{L, N} (z_1, z_2) (I^{\ast} \rho_x^L) (z_1) (I^{\ast} \rho_x^L) (z_2)
   \mathd z_1 \mathd z_2 - K^{L, N} \right] \mathbb{E} [\phi_0^L (x) \Phi
   (u_0)] \mathd x \]
and
\[ C^{L, N} = 2 \int_{\mathbb{T}} \varphi (x) \int_{\mathbb{T}^2} g_x^{L, N}
   (z_1, z_2) (I^{\ast} \rho_x^L) (z_1) \mathbb{E} [\phi_0^L (x) \mathD_{z_2}
   \Phi (u_0)] \mathd z_1 \mathd z_2 \mathd x . \]
Let
\[ K^{L, N} = \int_{\mathbb{T}^2} g_x^{L, N} (z_1, z_2) (I^{\ast} \rho_x^L)
   (z_1) (I^{\ast} \rho_x^L) (z_2) \mathd z_1 \mathd z_2 \]
to have $B^{L, N} = 0$. From the expression of $A^{L, N}$ and $C^{L, N}$ we
deduce that with this choice
\[ \mathbb{E} [r^{L, N} (u_0, \varphi)] = 0 \]
as required. By Cauchy--Schwarz we have
\[ (A^{L, N})^2 \leqslant \mathbb{E} \left\| \int_{\mathbb{T}} \varphi (x)
   (W_1 (g_x^{L, N} (z_1, \cdot)) \diamond \phi_0^L (x)) \mathd x
   \right\|_{H^{- 1}_{z_1}}^2 \mathbb{E} \| \mathD_{z_1} \Phi (u_0)
   \|_{H^1_{z_1}}^2 = (A_1^{L, N})^2 \| \Phi \|_1^2 \]
where
\[ (A_1^{L, N})^2 =\mathbb{E} \left\| \int_{\mathbb{T}} \varphi (x) W_1
   (g_x^{L, N} (z_1, \cdot)) \diamond \phi_0^L (x) \mathd x \right\|_{H^{-
   1}_{z_1}}^2 \]
\[ = \int_{\mathbb{T}^2} \varphi (x) \varphi (x') \mathbb{E} \left[ \langle
   \phi_0^L (x) \diamond W_1 (g_x^{L, N} (z_1, \cdot)), \phi_0^L (x') \diamond
   W_1 (g_{x'}^{L, N} (z_1, \cdot)) \rangle_{H^{- 1}_{z_1}} \right] \mathd x
   \mathd x' \]
and
\[ (C^{L, N})^2 \lesssim \mathbb{E} \left\| \int_{\mathbb{T}} \varphi (x)
   \phi^L_0 (x) \int_{\mathbb{T}} g_x^{L, N} (z_1, z_2) (I^{\ast} \rho_x^L)
   (z_1) \mathd z_1 \mathd x \right\|_{H^{- 1}_{z_2}}^2 \mathbb{E} \|
   \mathD_{z_2} \Phi (u_0) \|_{H^1_{z_2}}^2 \]
\[ \lesssim (C_1^{L, N})^2 \| \Phi \|_1^2 \]
with
\[ (C_1^{L, N})^2 =\mathbb{E} \left\| \int_{\mathbb{T}} \varphi (x) \phi^L_0
   (x) \int_{\mathbb{T}} g_x^{L, N} (z_1, z_2) (I^{\ast} \rho_x^L) (z_1)
   \mathd z_1 \mathd x \right\|_{H^{- 1}_{z_2}}^2 \]
So we have proven

\begin{lemma}
  We have
  \[ | \mathbb{E} [r^{L, N} (u_0, \varphi) \Phi (u_0)] | \lesssim \| \Phi
     \|_1 (A_1^{L, N} + C_1^{L, N}) \]
  and in particular
  \[ \| r^{L, N} (u_0, \varphi) \|_{- 1} \lesssim A_1^{L, N} + C_1^{L, N} . \]
\end{lemma}

Then it remains to bound each of these constants to show that they vanish in
the limit $N, L \rightarrow + \infty$ and to study the limit of $K^{L, N}$.
This is the aim of the following lemmas.

\begin{lemma}
  We have $\lim_{N \rightarrow + \infty} K^{L, N} = K^L$ and $\lim_{L
  \rightarrow + \infty} K^L = K = - 1 / 12$.
\end{lemma}

\begin{proof}
  Since
  \[ \int_{\mathbb{T}} (I^{\ast} \rho_x^L) (z_1) \rho^N_y (z_1) \mathd z_1 =
     \left\langle I^{\ast} \JJ^L_0 \delta_x, \JJ^N_0 \delta_y
     \right\rangle_{L^2} = \left\langle \delta_x, I \JJ^L_0  \JJ^N_0 \delta_y
     \right\rangle_{L^2} = \left( I \JJ^L_0  \JJ^N_0 \delta_y \right) (x) \]
  we have
  \[ K^{L, N} = \int_{\mathbb{T}^2} g_x^{L, N} (z_1, z_2) (I^{\ast} \rho_x^L)
     (z_1) (I^{\ast} \rho_x^L) (z_2) \mathd z_1 \mathd z_2 \]
  \[ = \int_{\mathbb{T}} \mathd y (\rho_x^L (y) \left( I \JJ^{L \wedge N}_0
     \delta_y \right) (x)^2 - \rho_x^{\infty} (y) \left( I \JJ^L_0 \delta_y
     \right) (x)^2) . \]
  Taking $N \rightarrow \infty$ we get
  \[ K^{L, N} \rightarrow K^L = \int_{\mathbb{T}} \mathd y (\rho_x^L (y)
     \left( I \JJ^L_0 \delta_y \right) (x)^2 - \rho_x^{\infty} (y) \left( I
     \JJ^L_0 \delta_y \right) (x)^2) . \]
  Now note that the integral of the first term in the expression of $K^L$
  vanishes, indeed:
  \[ \int_{\mathbb{T}} \mathd y \rho_x^L (y) \left( I \JJ^L_0 \delta_y
     \right) (x)^2 = \int_{\mathbb{T}} \mathd y \rho_x^L (y) \left( I \JJ^L_0
     \delta_x \right) (y)^2 \]
  \[ = \int_{\mathbb{T}} \mathd y \left( \partial_y I \JJ_0^L \delta_x
     \right) (y) \left( I \JJ^L_0 \delta_x \right) (y)^2 = \frac{1}{3}
     \int_{\mathbb{T}} \mathd y \partial_y \left( I \JJ^L_0 \delta_x \right)
     (y)^3 = 0. \]
  Moreover the second term satisfies
  \[ \int_{\mathbb{T}} \mathd y \rho_x^{\infty} (y) \left( I \JJ^L_0 \delta_y
     \right) (x)^2 = \left( I \JJ^L_0 \delta_x \right) (x)^2 - \frac{1}{2 \pi}
     \int_{\mathbb{T}} \mathd y \left( I \JJ^L_0 \delta_y \right) (x)^2 \]
  since $\rho_x^{\infty} = \delta_x - 1 / 2 \pi$. By symmetry also the $\left(
  I \JJ^L_0 \delta_x \right) (x)^2$ contribution vanishes, since
  \[ \left( I \JJ^L_0 \delta_x \right) (x) = \left( I^{\ast} \JJ^L_0 \delta_x
     \right) (x) = - \left( I^{} \JJ^L_0 \delta_x \right) (x) = 0. \]
  So we ends up with
  \[ K^L = - \frac{1}{2 \pi} \int_{\mathbb{T}} \mathd y \left( I \JJ^L_0
     \delta_x \right) (y)^2 = - \frac{1}{2 \pi} \left\langle I \JJ^L_0
     \delta_x, I \JJ^L_0 \delta_x \right\rangle = - \frac{1}{4 \pi^2} \sum_{0
     < | k |} \frac{| \hat{\rho}^L (k) |^2}{| k |^2} \]
  since
  \[ \FF \left( I \JJ^L_0 \delta_x \right) (k) = \frac{\hat{\rho}^L (k)}{ik}
     e_k (x)^{\ast} . \]
  And as $L \rightarrow \infty$ we obtain $\hat{\rho}^L (k) \rightarrow 1$ for
  all $k \neq 0$ so
  \[ K^L \rightarrow - \frac{1}{2 \pi^2} \sum_{n = 1}^{\infty} \frac{1}{n^2} =
     - \frac{1}{12} = K. \]
  
\end{proof}

\begin{lemma}
  We have $A_1^{L, N} \rightarrow 0$ as $N \rightarrow \infty$ and $L
  \rightarrow \infty$.
\end{lemma}

\begin{proof}
  
  \[ (A_1^{L, N})^2 =\mathbb{E} \left\| \int_{\mathbb{T}} \varphi (x) W_1
     (g_x^{L, N} (z_1, \cdot)) \diamond \phi_0^L (x) \mathd x \right\|_{H^{-
     1}_{z_1}}^2 \]
  \[ = \int_{\mathbb{T}^2} \varphi (x) \varphi (x') \mathbb{E} \left[
     \langle \phi_0^L (x) \diamond W_1 (g_x^{L, N} (z_1, \cdot)), \phi_0^L
     (x') \diamond W_1 (g_{x'}^{L, N} (z_1, \cdot)) \rangle_{H^{- 1}_{z_1}}
     \right] \mathd x \mathd x' \]
  Integrating by parts the $W_1$ terms and taking into account the
  cancellations due to the partial $\diamond$ contractions we get
  \[ \mathbb{E} \left[ \langle \phi_0^L (x) \diamond W_1 (g_x^{L, N} (z_1,
     \cdot)), \phi_0^L (x') \diamond W_1 (g_{x'}^{L, N} (z_1, \cdot))
     \rangle_{H^{- 1}_{z_1}} \right] \]
  \[ =\mathbb{E} [\phi_0^L (x) \phi_0^L (x')] \int_{\mathbb{T}^2} \langle
     g_x^{L, N} (z_1, z_2), g_{x'}^{L, N} (z_1, z_2) \rangle_{H^{- 1}_{z_1}}
     \mathd z_2 \]
  \[ + \int_{\mathbb{T}^2} \mathbb{E} [(\mathD_{z_3} \phi_0^L (x))
     (\mathD_{z_2} \phi_0^L (x'))] \langle g_x^{L, N} (z_1, z_2), g_{x'}^{L,
     N} (z_1, z_3) \rangle_{H^{- 1}_{z_1}} \mathd z_2 \mathd z_3 \]
  The second term can be written as
  \[ \mathbb{E} [\phi_0^L (x) \phi_0^L (x')] \int_{\mathbb{T}^2} (I^{\ast}
     \rho_x^L) (z_3) (I^{\ast} \rho_{x'}^L) (z_2) \langle g_x^{L, N} (z_1,
     z_2), g_{x'}^{L, N} (z_1, z_3) \rangle_{H^{- 1}_{z_1}} \mathd z_2 \mathd
     z_3 \]
  so letting
  \[ V^L (x, x') = \varphi (x) \varphi (x') \mathbb{E} [_{} \phi_0^L (x)
     \phi_0^L (x')] \]
  we have
  \[ (A_1^{L, N})^2 = \int_{\mathbb{T}^2} V^L (x, x') \int_{\mathbb{T}^2}
     (I^{\ast} \rho_x^L) (z_3) (I^{\ast} \rho_{x'}^L) (z_2) \langle g_x^{L, N}
     (z_1, z_2), g_{x'}^{L, N} (z_1, z_3) \rangle_{H^{- 1}_{z_1}} \mathd z_2
     \mathd z_3 \mathd x \mathd x' \]
  \[ + \int_{\mathbb{T}^2} V^L (x, x') \int_{\mathbb{T}^2} \langle g_x^{L,
     N} (z_1, z_2), g_{x'}^{L, N} (z_1, z_2) \rangle_{H^{- 1}_{z_1}} \mathd
     z_2 \mathd x \mathd x' 
   = A_{1, 1} + A_{1, 2} \]
  Let us consider $A_{1, 1}$ first :
  \[ A_{1, 1} = \int_{\mathbb{T}^4} V^L (x, x') \langle G_{x, x'}, G_{x', x}
     \rangle_{H^{- 1}} \mathd x \mathd x' \]
  where
  \[ G_{x, x'} (z_1) = \int_{\mathbb{T}} (I^{\ast} \rho_{x'}^L) (z_2) g_x^{L,
     N} (z_1, z_2) \mathd z_2 . \]
  Now observe that
  \[ G_{x, x'} (z_1) = \int_{\mathbb{T}} \mathd y \rho_x^L (y) \rho^N_y (z_1)
     (\JJ^N_0 I^{\ast} \rho_{x'}^L) (y) - \int_{\mathbb{T}} \mathd y
     \rho_x^{\infty} (y) \rho^L_y (z_1) (I^{\ast} \rho_{x'}^L) (y) \]
  so the $H^{- 1}$ scalar product in $A_{1, 1}$ can be expanded as
  \[ \langle G_{x, x'}, G_{x', x} \rangle_{H^{- 1}} = \int_{\mathbb{T}}
     \mathd y \rho_x^L (y) (\JJ^N_0 I^{\ast} \rho_{x'}^L) (y)
     \int_{\mathbb{T}} \mathd y' \rho_{x'}^L (y') (\JJ^N_0 I^{\ast} \rho_x^L)
     (y') \langle \rho^N_y (z_1), \rho^N_{y'} (z_1) \rangle_{H^{- 1}_{z_1}} \]
  \[ - \int_{\mathbb{T}} \mathd y \rho_x^{\infty} (y) (I^{\ast} \rho_{x'}^L)
     (y) \int_{\mathbb{T}} \mathd y' \rho_{x'}^L (y') (\JJ^N_0 I^{\ast}
     \rho_x^L) (y') \langle \rho^L_y (z_1), \rho^N_{y'} (z_1) \rangle_{H^{-
     1}_{z_1}} \]
  \[ - \int_{\mathbb{T}} \mathd y \rho_x^L (y) (\JJ^N_0 I^{\ast} \rho_{x'}^L)
     (y) \int_{\mathbb{T}} \mathd y' \rho_{x'}^{\infty} (y') (I^{\ast}
     \rho_x^L) (y') \langle \rho^N_y (z_1), \rho^L_{y'} (z_1) \rangle_{H^{-
     1}_{z_1}} \]
  \[ + \int_{\mathbb{T}} \mathd y \rho_x^{\infty} (y) (I^{\ast} \rho_{x'}^L)
     (y) \int_{\mathbb{T}} \mathd y' \rho_{x'}^{\infty} (y') (I^{\ast}
     \rho_x^L) (y') \langle \rho^L_y (z_1), \rho^L_{y'} (z_1) \rangle_{H^{-
     1}_{z_1}} \]
  Taking $N \rightarrow \infty$ we get
  \[ \langle G_{x, x'}, G_{x', x} \rangle_{H^{- 1}} \rightarrow
     \int_{\mathbb{T}} \mathd y \rho_x^L (y) (I^{\ast} \rho_{x'}^L) (y)
     \int_{\mathbb{T}} \mathd y' \rho_{x'}^L (y') (I^{\ast} \rho_x^L) (y')
     \langle \rho^{\infty}_y (z_1), \rho^{\infty}_{y'} (z_1) \rangle_{H^{-
     1}_{z_1}} \]
  \[ - \int_{\mathbb{T}} \mathd y \rho_x^{\infty} (y) (I^{\ast} \rho_{x'}^L)
     (y) \int_{\mathbb{T}} \mathd y' \rho_{x'}^L (y') (I^{\ast} \rho_x^L)
     (y') \langle \rho^L_y (z_1), \rho^{\infty}_{y'} (z_1) \rangle_{H^{-
     1}_{z_1}} \]
  \[ - \int_{\mathbb{T}} \mathd y \rho_x^L (y) (I^{\ast} \rho_{x'}^L) (y)
     \int_{\mathbb{T}} \mathd y' \rho_{x'}^{\infty} (y') (I^{\ast} \rho_x^L)
     (y') \langle \rho^{\infty}_y (z_1), \rho^L_{y'} (z_1) \rangle_{H^{-
     1}_{z_1}} \]
  \[ + \int_{\mathbb{T}} \mathd y \rho_x^{\infty} (y) (I^{\ast} \rho_{x'}^L)
     (y) \int_{\mathbb{T}} \mathd y' \rho_{x'}^{\infty} (y') (I^{\ast}
     \rho_x^L) (y') \langle \rho^L_y (z_1), \rho^L_{y'} (z_1) \rangle_{H^{-
     1}_{z_1}} \]
  \[ = \left\langle I \JJ_0 (\rho^L_x (I^{\ast} \rho_{x'}^L)), I \JJ_0
     (\rho_{x'}^L (I^{\ast} \rho_x^L)) \right\rangle - \left\langle I \JJ_0^L
     (\rho^{\infty}_x (I^{\ast} \rho_{x'}^L)), I \JJ_0 (\rho_{x'}^L (I^{\ast}
     \rho_x^L)) \right\rangle \]
  \[ - \left\langle I \JJ_0 (\rho^L_x (I^{\ast} \rho_{x'}^L)), I \JJ_0^L
     (\rho_{x'}^{\infty} (I^{\ast} \rho_x^L)) \right\rangle + \left\langle I
     \JJ_0^L (\rho^{\infty}_x (I^{\ast} \rho_{x'}^L)), I \JJ_0^L
     (\rho_{x'}^{\infty} (I^{\ast} \rho_x^L)) \right\rangle . \]
  The point is to show that, as $L \rightarrow + \infty$,
  \[ \left\langle I \JJ_0 (\rho^L_x (I^{\ast} \rho_{x'}^L)), I \JJ_0
     (\rho_{x'}^L (I^{\ast} \rho_x^L)) \right\rangle \rightarrow \left\langle
     I \JJ_0 (\rho^{\infty}_x (I^{\ast} \rho_{x'}^{\infty})), I \JJ_0
     (\rho_{x'}^{\infty} (I^{\ast} \rho_x^{\infty})) \right\rangle \]
  and the same limit for the other three quantities. Writing things in Fourier
  space we have
  \[ J_L (x - x') = \left\langle I \JJ_0 (\rho^L_x (I^{\ast} \rho_{x'}^L)), I
     \JJ_0 (\rho_{x'}^L (I^{\ast} \rho_x^L)) \right\rangle \]
  \[ = \frac{1}{2 \pi} \sum_{\tmscript{\begin{array}{c}
       0 < | k_1 |, | k_2 |, | k_3 |, | k_4 |\\
       k_1 + k_2 + k_3 + k_4 = 0
     \end{array}}} \frac{\mathbb{I}_{k_1 + k_2 \neq 0}}{(k_1 + k_2)^2}
     e^{ik_1 x} \frac{e^{ik_2 x'}}{- ik_2} e^{ik_3 x'} \frac{e^{ik_4 x}}{-
     ik_4} \hat{\rho}^L (k_1) \hat{\rho}^L (k_2) \hat{\rho}^L (k_3)
     \hat{\rho}^L (k_4) \]
  \[ = \frac{1}{2 \pi} \sum_{\tmscript{\begin{array}{c}
       0 < | k_1 |, | k_2 |, | k_3 |, | k_4 |\\
       k_1 + k_2 + k_3 + k_4 = 0
     \end{array}}} \frac{\mathbb{I}_{k_1 + k_2 \neq 0}}{(k_1 + k_2)^2}
     \frac{e^{i (k_2 + k_3) (x' - x)}}{- ik_2} \frac{1}{- ik_4} \hat{\rho}^L
     (k_1) \hat{\rho}^L (k_2) \hat{\rho}^L (k_3) \hat{\rho}^L (k_4) \]
  Now taking Littlewood--Paley blocks $(\Delta_q)_{q \geqslant - 1}$ of this
  quantity we get
  \[ \Delta_q J_L (x) = \frac{1}{2 \pi} \sum_{\tmscript{\begin{array}{c}
       0 < | k_1 |, | k_2 |, | k_3 |, | k_4 |\\
       k_1 + k_2 + k_3 + k_4 = 0
     \end{array}}} \frac{\mathbb{I}_{k_1 + k_2 \neq 0}}{(k_1 + k_2)^2}
     \frac{e^{i (k_2 + k_3) (x' - x)}}{- ik_2} \frac{\theta_q (k_2 + k_3)}{-
     ik_4} \hat{\rho}^L (k_1) \hat{\rho}^L (k_2) \hat{\rho}^L (k_3)
     \hat{\rho}^L (k_4) . \]
  Where $\theta_q$ is the Fourier multiplier associated to $\Delta_q$ by
  $\Delta_q = \theta_q (\mathD)$. Taking into accout that $\sup_k |
  \hat{\rho}^L (k) | < + \infty$ if we show that the quantity
  \[ S_q = \frac{1}{2 \pi} \sum_{\tmscript{\begin{array}{c}
       0 < | k_1 |, | k_2 |, | k_3 |, | k_4 |\\
       k_1 + k_2 + k_3 + k_4 = 0
     \end{array}}} \mathbb{I}_{k_1 + k_2 \neq 0} \frac{\theta_q (k_2 +
     k_3)}{(k_1 + k_2)^2 | k_2 |  | k_3 |} \]
  is finite then we can conclude by dominated convergence. Note that
  \[ S_q \lesssim 2^{q \varepsilon} \sum_{0 < | k_1 |, | k_2 |, | k_3 |}
     \frac{\mathbb{I}_{k_1 + k_2 \neq 0, k_2 + k_3 \neq 0}}{(k_1 + k_2)^2 |
     k_2 |  | k_3 | | k_2 + k_3 |^{\varepsilon}} \]
  for any $\varepsilon > 0$ since
  \[ | \theta_q (k) | = | \theta (2^{- q} k) | \lesssim (1 + | 2^{- q} k |)^{-
     \varepsilon} \lesssim 2^{q \varepsilon} | k |^{- \varepsilon} \]
  due to the fact that the function $\theta$ is supported in a ball of finite
  radius. We can now perform the sum over $k_3$ and get
  \[ S_q \lesssim 2^{q \varepsilon} \sum_{0 < | k_1 |, | k_2 |}
     \frac{\mathbb{I}_{k_1 + k_2 \neq 0}}{(k_1 + k_2)^2 | k_2 | } \]
  since
  \[ \sum_{0 < | k_3 |} \frac{\mathbb{I}_{k_2 + k_3 \neq 0}}{| k_3 | | k_2 +
     k_3 |^{\varepsilon}} \lesssim \frac{1}{| k_2 |^{\varepsilon}} \]
  uniformly in $k_2 \neq 0$ and $L$. Now performing the other two sums we show
  simply that
  \[ S_q \lesssim 2^{q \varepsilon} \sum_{0 < | k_2 |} \frac{1}{| k_2 |^{1 +
     \varepsilon} } \lesssim 2^{q \varepsilon} \]
  where we used the fact that
  \[ \sum_{0 < | k_1 |} \frac{\mathbb{I}_{k_1 + k_2 \neq 0}}{(k_1 + k_2)^2 }
     \leqslant \sum_{k_1 \in \mathbb{Z}} \frac{\mathbb{I}_{k_1 + k_2 \neq
     0}}{(k_1 + k_2)^2 } = \sum_{k_1 \in \mathbb{Z}} \frac{\mathbb{I}_{k_1
     \neq 0}}{(k_1)^2 } \lesssim 1. \]
  By dominated convergence we can conclude that
  \[ \left\langle I \JJ_0 (\rho^L_x (I^{\ast} \rho_{x'}^L)), I \JJ_0
     (\rho_{x'}^L (I^{\ast} \rho_x^L)) \right\rangle \rightarrow \left\langle
     I \JJ_0 (\rho^{\infty}_x (I^{\ast} \rho_{x'}^{\infty})), I \JJ_0
     (\rho_{x'}^{\infty} (I^{\ast} \rho_x^{\infty})) \right\rangle \]
  as distributions of arbitrarily small negative order of the variable $x -
  x'$. This is ok since we test this quantity against the function $V^L$ which
  belongs to $C^{1 / 2 -} (\mathbb{T}^2)$ uniformly in $L$. The other terms
  can be handled similarly and thus we conclude that $A_{1, 1} \rightarrow 0$.
  
  \
  
  Let us turn now to $A_{1, 2}$. The relevant computation is the following:
  \[ \int_{\mathbb{T}^2} \langle g_x^{L, N} (z_1, z_2), g_{x'}^{L, N} (z_1,
     z_2) \rangle_{H^{- 1}_{z_1}} \mathd z_2 = \]
  \[ \int_{\mathbb{T}} \mathd y \rho_x^L (y) \int_{\mathbb{T}} \mathd y'
     \rho_{x'}^L (y') \langle \rho^N_y (z_1), \rho^N_{y'} (z_1) \rangle_{H^{-
     1}_{z_1}} \langle \rho^N_y (z_2), \rho^N_{y'} (z_2) \rangle_{L^2_{z_2}}
  \]
  \[ - \int_{\mathbb{T}} \mathd y \rho_x^L (y) \int_{\mathbb{T}} \mathd y'
     \rho_{x'}^{\infty} (y') \langle \rho^N_y (z_1), \rho^L_{y'} (z_1)
     \rangle_{H^{- 1}_{z_1}} \langle \rho^N_y (z_2), \rho^L_{y'} (z_2)
     \rangle_{L^2_{z_2}} \]
  \[ - \int_{\mathbb{T}} \mathd y \rho_x^{\infty} (y) \int_{\mathbb{T}}
     \mathd y' \rho_{x'}^L (y') \langle \rho^L_y (z_1), \rho^N_{y'} (z_1)
     \rangle_{H^{- 1}_{z_1}} \langle \rho^L_y (z_2), \rho^N_{y'} (z_2)
     \rangle_{L^2_{z_2}} \]
  \[ + \int_{\mathbb{T}} \mathd y \rho_x^{\infty} (y) \int_{\mathbb{T}}
     \mathd y' \rho_{x'}^{\infty} (y') \langle \rho^L_y (z_1), \rho^L_{y'}
     (z_1) \rangle_{H^{- 1}_{z_1}} \langle \rho^L_y (z_2), \rho^L_{y'} (z_2)
     \rangle_{L^2_{z_2}} \]
  Now consider one of the terms:
  \[ \int_{\mathbb{T}} \mathd y \rho_x^L (y) \int_{\mathbb{T}} \mathd y'
     \rho_{x'}^L (y') \langle \rho^N_y (z_1), \rho^N_{y'} (z_1) \rangle_{H^{-
     1}_{z_1}} \langle \rho^N_y (z_2), \rho^N_{y'} (z_2) \rangle_{L^2_{z_2}}
  \]
  \[ = \int_{\mathbb{T}} \mathd y \rho_x^L (y) \int_{\mathbb{T}} \mathd y'
     \rho_{x'}^L (y') \rho^N_y (y') (I^{\ast} I \rho^N_y) (y') \]
  since
  \[ \langle \rho^N_y (z_1), \rho^N_{y'} (z_1) \rangle_{H^{- 1}_{z_1}} =
     \langle I \rho^N_y, I \rho^N_{y'} \rangle_{L^2} = \langle I^{\ast} I
     \rho^N_y, \delta_{y'} \rangle_{L^2} = (I^{\ast} I \rho^N_y) (y') \]
  As $N \rightarrow \infty$ this quantity becomes
  \[ \rightarrow \int_{\mathbb{T}} \mathd y \rho_x^L (y) \int_{\mathbb{T}}
     \mathd y' \rho_{x'}^L (y') \rho^{\infty}_y (y') (I^{\ast} I
     \rho^{\infty}_y) (y') \]
  \[ = \int_{\mathbb{T}} \mathd y \rho_x^L (y) \rho_{x'}^L (y) (I^{\ast} I
     \rho^{\infty}_y) (y) - \frac{1}{2 \pi} \int_{\mathbb{T}} \mathd y
     \rho_x^L (y) \int_{\mathbb{T}} \mathd y' \rho_{x'}^L (y') (I^{\ast} I
     \rho^{\infty}_y) (y') \]
  \[ = (I^{\ast} I \rho^{\infty}_0) (0) \int_{\mathbb{T}} \mathd y \rho_x^L
     (y) \rho_{x'}^L (y) - \frac{1}{2 \pi} (I^{\ast} I \rho^L_x) (x') \]
  \[ = (I^{\ast} I \rho^{\infty}_0) (0) \rho_x^L (x') - \frac{1}{2 \pi}
     (I^{\ast} I \rho^L_x) (x') \]
  since $(I^{\ast} I \rho^{\infty}_y) (y) = (I^{\ast} I \rho^{\infty}_0) (0)$.
  As $L \rightarrow \infty$ we obtain
  \[ = (I^{\ast} I \rho^{\infty}_0) (0) \rho_x^{\infty} (x') - \frac{1}{2 \pi}
     (I^{\ast} I \rho^{\infty}_x) (x') \]
  For the other terms the discussion of the convergence is similar and the
  limit is the same so we obtain that
  \[ \int_{\mathbb{T}^2} \langle g_x^{L, N} (z_1, z_2), g_{x'}^{L, N} (z_1,
     z_2) \rangle_{H^{- 1}_{z_1}} \mathd z_2 \rightarrow 0 \]
  as a distribution in the $x, x'$ variables and we conclude that $A_{1, 2}
  \rightarrow 0$.
\end{proof}

\begin{lemma}
  We have $C_1^{L, N} \rightarrow 0$ as $N \rightarrow \infty$ and $L
  \rightarrow \infty$.
\end{lemma}

\begin{proof}
  Let
  \[ G_x (z_2) = \int_{\mathbb{T}} g_x^{L, N} (z_1, z_2) (I^{\ast} \rho_x^L)
     (z_1) \mathd z_1 = G_{x, x} (z_2) \]
  and recall that
  \[ (C_1^{L, N})^2 =\mathbb{E} \left\| \int_{\mathbb{T}} \varphi (x)
     \phi^L_0 (x) G_x (z_2) \mathd x \right\|_{H^{- 1}_{z_2}}^2 =
     \int_{\mathbb{T}^2} V^L (x, x') \langle G_x, G_{x'} \rangle_{H^{- 1}}
     \mathd x \mathd x' \]
  In order to estimate the scalar product can take $f \in H^1$ and consider
  $\int G_{x, x} (z) f (z) \mathd z$ instead. Recall that
  
  
  \[ \int G_{x, x} (z) f (z) \mathd z = \int_{\mathbb{T}} \mathd y \rho_x^L
     (y) \JJ^N_0 f (y) (\JJ^N_0 I^{\ast} \rho_x^L) (y) - \int_{\mathbb{T}}
     \mathd y \rho_x^{\infty} (y) \JJ^L_0 f (y) (I^{\ast} \rho_x^L) (y) \]
  replace $\rho^{\infty}_x (y)$ by $\delta_x (y) - 1 / 2 \pi$ and $\rho^L_x
  (y)$ by $\delta_x^L (y) - 1 / 2 \pi$ to get
  \[ = \int_{\mathbb{T}} \mathd y \delta_x^L (y) \left( \JJ^N_0 f (y) \right)
     (\JJ^N_0 I^{\ast} \rho_x^L) (y) - \left( \JJ^L_0 f (x) \right) (I^{\ast}
     \rho_x^L) (x) \]
  \[ - \frac{1}{2 \pi} \int_{\mathbb{T}} \mathd y \JJ^N_0 f (y) (\JJ^N_0
     I^{\ast} \rho_x^L) (y) + \frac{1}{2 \pi} \int_{\mathbb{T}} \mathd y
     \JJ^L_0 f (y) (I^{\ast} \rho_x^L) (y) \]
  \[ = \int_{\mathbb{T}} \mathd y \delta_x^L (y) \left( \JJ^N_0 f (y) -
     \JJ^N_0 f (x) \right) (\JJ^N_0 I^{\ast} \rho_x^L) (y) \]
  \[ + \left( \JJ^N_0 f (x) \right) (I^{\ast} \left( \delta^L \ast \JJ^N_0
     \rho_x^L \right)) (x) - \left( \JJ^L_0 f (x) \right) (I^{\ast} \rho_x^L)
     (x) \]
  \[ - \frac{1}{2 \pi} \int_{\mathbb{T}} \mathd y \JJ^N_0 f (y) (\JJ^N_0
     I^{\ast} \rho_x^L) (y) + \frac{1}{2 \pi} \int_{\mathbb{T}} \mathd y
     \JJ^L_0 f (y) (I^{\ast} \rho_x^L) (y) . \]
  We note that, by symmetry,
  \[ \left( \JJ^N_0 f (x) \right) (I^{\ast} \left( \delta^L \ast \JJ^N_0
     \rho_x^L \right)) (x) = \left( \JJ^L_0 f (x) \right) (I^{\ast} \rho_x^L)
     (x) = 0 \]
  and moreover that in this expression we have that
  \[ \sup_x \left| \left\langle \JJ^L_0 f, I^{\ast} \rho_x^L \right\rangle -
     \langle f, \Theta_x^{} \rangle \right| + \sup_x \left| \left\langle
     \JJ^N_0 f, \JJ^N_0 I^{\ast} \rho_x^L \right\rangle - \langle f,
     \Theta_x^{} \rangle \right| \lesssim \| f \|_{H^1} o (1)
  \]
 as $L\wedge N \to \infty$. Finally it remains to bound
  \[ H = \int_{\mathbb{T}} \mathd y \delta_x^L (y) \left( \JJ^N_0 f (y) -
     \JJ^N_0 f (x) \right) (\JJ^N_0 I^{\ast} \rho_x^L) (y) \]
  However using $H^1 \subseteq C^{1 / 2}$ compactly and that $\left\| \JJ^N_0
  I^{\ast} \rho^L_x \right\|_{\infty} < + \infty$ uniformy in $x, L, N$ we
  have
  \[ | H | \lesssim \left\| \JJ^N_0 f \right\|_{H^1} \int_{\mathbb{T}} \mathd
     y \delta_x^L (y) | y - x |^{1 / 2} \lesssim \| f \|_{H^1} L^{- 1 / 2}
     \rightarrow 0. \]
  In conclusion
  \[ \sup_{x \in \mathbb{T}} \left| \int G_{x, x'} (z) f (z) \mathd z
     \right|_{} \lesssim \| f \|_{H^1} o (1) \]
  which implies
  $\sup_{x \in \mathbb{T}} \langle G_{x, x}, G_{x, x} \rangle_{H^{- 1}} = o
     (1) $
  and thus we can conclude that $C_1^{L, N} \rightarrow 0$.
\end{proof}

\begin{thebibliography}{10}
  \bibitem[1]{BertiniGiacomin1997}Lorenzo~Bertini  and 
  Giambattista~Giacomin.{\newblock} Stochastic Burgers and KPZ equations from
  particle systems.{\newblock} \tmtextit{Comm. Math. Phys.}, 183(3):571--607,
  1997.{\newblock}
  
  \bibitem[2]{funaki_kpz_2014}Tadahisa~Funaki  and  Jeremy~Quastel.{\newblock}
  KPZ equation, its renormalization and invariant measures.{\newblock}
  \tmtextit{ArXiv:1407.7310 [math]}, , jul 2014.{\newblock} ArXiv:
  1407.7310.{\newblock}
  
  \bibitem[3]{goncalves_nonlinear_2013}Patr$\acute{\text{{\i}}}$cia~Gon{\c
  c}alves  and  Milton~Jara.{\newblock} Nonlinear Fluctuations of Weakly
  Asymmetric Interacting Particle Systems.{\newblock} \tmtextit{Archive for
  Rational Mechanics and Analysis}, 212(2):597--644, dec 2013.{\newblock}
  
  \bibitem[4]{gubinelli_regularization_2013}M.~Gubinelli  and 
  M.~Jara.{\newblock} Regularization by noise and stochastic Burgers
  equations.{\newblock} \tmtextit{Stochastic Partial Differential Equations:
  Analysis and Computations}, 1(2):325--350, jun 2013.{\newblock}
  
  \bibitem[5]{gubinelli_lectures_2015}M.~Gubinelli  and 
  N.~Perkowski.{\newblock} Lectures on singular stochastic PDEs.{\newblock}
  \tmtextit{ArXiv:1502.00157 [math]}, , jan 2015.{\newblock} ArXiv:
  1502.00157.{\newblock}
  
  \bibitem[6]{gubinelli_paracontrolled_2015}Massimiliano~Gubinelli,
  Peter~Imkeller  and  Nicolas~Perkowski.{\newblock} Paracontrolled
  distributions and singular PDEs.{\newblock} \tmtextit{Forum of Mathematics,
  Pi}, 3, 2015.{\newblock}
  
  \bibitem[7]{hairer_solving_2011}Martin~Hairer.{\newblock} Solving the KPZ
  equation.{\newblock} \tmtextit{Ann. Math.}, 178(2):559--664,
  2013.{\newblock}
  
  \bibitem[8]{Hairer2014Regularity}Martin~Hairer.{\newblock} A theory of
  regularity structures.{\newblock} \tmtextit{Invent. Math.}, 198(2):269--504,
  2014.{\newblock}
  
  \bibitem[9]{hairer-quastel}Martin~Hairer  and  Jeremy~Quastel.{\newblock}
  Weak universality of the kpz equation.{\newblock} Announcement,
  2014.{\newblock}
  
  \bibitem[10]{prato_stochastic_2008}Guiseppe Da~Prato  and 
  Jerzy~Zabczyk.{\newblock} \tmtextit{Stochastic Equations in Infinite
  Dimensions}.{\newblock} Cambridge University Press, feb 2008.{\newblock}
  
  \bibitem[11]{russo_elements_2007}Francesco~Russo  and 
  Pierre~Vallois.{\newblock} Elements of Stochastic Calculus via
  Regularization.{\newblock} In  Catherine~Donati-Martin, Michel~{\'E}mery,
  Alain~Rouault  and  Christophe~Stricker , editors , \tmtextit{S{\'e}minaire
  de Probabilit{\'e}s XL},  volume  1899  of \tmtextit{Lecture Notes in
  Mathematics},  pages  147--185. Springer Berlin / Heidelberg,
  2007.{\newblock}
  
  \bibitem[12]{fluctuations_2012}\tmtextit{Fluctuations in Markov Processes:
  Time Symmetry and Martingale Approximation}.{\newblock} Springer, Heidelberg
  Germany ; New York, 2012 edition  edition , jul 2012.{\newblock}
\end{thebibliography}

\end{document}
